Until the conclusion of this work, we concern ourselves with the airfoil skeleton theory of \textsc{Birnbaum}, wherein the aerodynamical calculations are performed on the mean camber line only.
For it shall be shown that the matter of flow due to thickness and inclination and camber are fundamentally different, and indeed separable.
Consider an airfoil given as a closed curve in the complex plane, whose chord coïncides with the abscissal coördinate.
The mean camber line is the curve that for any abscissal value lies exactly midway between the two ordinal values of the airfoil.
It is evident that this curve describing the airfoil can then be depicted entirely by the sum of the mean camber line and the thickness function.

In most practical scenarios, one would not be given a \textsc{Žukovskij} transform with which to analytically determine the mean camber line.
One would instead be given an airfoil in terms of discrete points along its contour.
The crucial observation to be made is that these points serve as vertices for a polygon inscribed in the airfoil.
If one were given a distribution of vertices such that the number along the suction side and pressure side were equal, one could consider the midpoint along the diagonal of corresponding vertices, numbered from leading edge to trailing edge, or \emph{vice versa}, as proposed by \textsc{Hess}.%
\footnote{%
        \cite{hess1974problem} \textsc{Hess} (1974), see fig.7a
}
Connecting each such midpoint from leading edge to trailing edge would yield an acceptable approximation to the mean camber line.
Should the number of provided vertices be different on the suction and pressure sides, however, the procedure through which to obtain the mean camber line is not as evident.
We shall determine an approximation to the mean camber line through studying the medial axes of polygons.
Further, we conjecture these to be the optimal approximations, for the error of the medial axis to the mean camber line cannot be larger than the error of the inscribed polygon relative to the true curve of the airfoil.
The medial axis was conceived of by \textsc{Blum} in a selcouth passage of writing on biological and visual shape description.%
\footnote{%
        \cite{blum1967transformation} \textsc{Blum} (1967)
}
Its definition is as follows.%
\footnote{%
        \cite{li1982medial} \textsc{Li} (1982), 
}

\begin{definition}[Medial axis]
  Given a simple polygon $\Omega^{\dag}$, its medial axis is the set of points $\{\xvec_i\}$ internal to it, such that there are at least two points on its boundary that are equdistant from and are closest to $\{ \xvec_i \}$.
\end{definition}

\noindent
The peculiar choice of notation of the polygon stems from the interior of the wing, whose cross section is the airfoil at hand, being the complement of the manifold inhabiting the fluid of interest.
Forcing a simple polygon, we employ the \textsc{Jordan} curve theorem,%
\footnote{%
        \cite{goluzin1969geometric} \textsc{Goluzin}, pp.8\&36
}
yielding a consistent notion of interior.
We can then implement an algorithm to determine the \textsc{Voronoj} diagram, whose set of edges is a superset of the medial axis,%
\footnote{%
        \cite{li1982medial} \textsc{Li} (1982), Corollary 3.
}
simply terminating branches thereof incident with reflex vertices.

\begin{conjecture}
        For an airfoil sampled at a finite number of points, the medial axis of the polygon thereof is the best approximation of its mean camber line.
\end{conjecture}

\begin{figure}[h]
        \centering
        \scalebox{1}{\input{airfoil_panel.tikz}}
        \caption{Airfoil}
\end{figure}
